\chapter{Mortgage Modelling}\label{ch:rc:mortgage-modelling}

On 13 June 2003 the ten-year Treasury yielded 3.13\%; on 2 September, 4.61\%. Part
of that move was the mortgage market hedging itself. American homeowners can repay
their fixed-rate mortgages at any time; when rates fall they refinance, when rates
rise they stay, so a mortgage-backed security shortens when rates fall and lengthens
when they rise. Holders who keep their portfolios' duration fixed must receive fixed
as rates fall and pay fixed as they rise, and in the summer of 2003 they paid into a
market that was already selling. A Federal Reserve study that year found that swap
volatility rose when prepayment risk was high, consistent with hedging amplifying
rate moves for months at a time. One Quant Book 2, chapter 12, described the market,
its prepayment conventions and its negative convexity with a single rate path. This
chapter models prepayment on many paths, prices a pool by option-adjusted spread,
splits it into its interest and principal, and measures the hedge a 100-basis-point
rise demands.

\section{Anatomy of a prepayment model}

\begin{definition}[Prepayment model]\label{def:rc:mortgage-modelling:model}
A \emph{prepayment model}\index{prepayment model} gives the conditional prepayment
rate of a pool each month as a function of the pool's characteristics (age, rate,
loan size) and of the rate environment on the path; it is estimated on historical
loan-level data and run path by path in valuation.
\end{definition}

\begin{definition}[Housing turnover, refinancing incentive]\label{def:rc:mortgage-modelling:turnover}
\emph{Housing turnover}\index{housing turnover} is the part of prepayment that comes
from home sales, defaults and curtailments, rising with loan age to a plateau and
insensitive to rates. The \emph{refinancing incentive}\index{refinancing incentive} is
the pool's weighted-average coupon (WAC) minus the mortgage rate a borrower could get
today: refinancing is worth it when the incentive exceeds the costs of doing so.
\end{definition}

\begin{definition}[Prepayment S-curve, burnout]\label{def:rc:mortgage-modelling:scurve}
The \emph{prepayment S-curve}\index{prepayment S-curve} is the refinancing speed as a
function of the incentive: near zero out of the money, rising steeply over a few tens
of basis points of incentive, flattening at a maximum. \emph{Burnout}\index{burnout}
is the decline of a pool's response to a given incentive after it has already been
exposed to incentives: the borrowers who could and would refinance have left, and
those who remain are slower.
\end{definition}

The chapter's model adds a turnover part at the PSA ramp (One Quant Book 2, chapter
12) to a refinancing S-curve with a maximum of 45\% CPR and a centre at 75 basis
points of incentive, damped by \omterm{def:rc:mortgage-modelling:scurve}{burnout} $e^{-0.25\,c}$ with $c$ the cumulative positive
incentive in percentage-point years
(\cref{fig:rc:mortgage-modelling:scurve}).

\begin{definition}[Current coupon, primary--secondary spread]\label{def:rc:mortgage-modelling:cc}
The \emph{current coupon}\index{current coupon} is the yield of a to-be-announced
mortgage security priced at par, interpolated between the coupons that trade. The
\emph{primary--secondary spread}\index{primary--secondary spread} is the mortgage rate
offered to borrowers (the primary market) minus the current coupon (the secondary
market): the lenders' and guarantors' margin, which widens when origination capacity
is scarce (in 2020 it ran 73 to 81 basis points above what demand explained, from March to September). A \omterm{def:rc:mortgage-modelling:model}{prepayment model} maps the rate paths to borrowers' mortgage rates
through both; the chapter uses the ten-year par swap rate plus a fixed 175 basis
points.
\end{definition}

\omcode{code/firm/mbsoas/firm_mbsoas.py}{36}{49}{The prepayment model: turnover plus a refinancing S-curve damped by burnout.}\label{lst:rc:mortgage-modelling:model}

\begin{omfigure}
\begin{tikzpicture}
\begin{axis}[width=11cm,height=5cm,
  xlabel={refinancing incentive: WAC minus mortgage rate (percentage points)},ylabel={refinancing CPR (\%)},
  tick label style={font=\small},label style={font=\small},
  xmin=-2,xmax=3,ymin=0,ymax=50,grid=major,grid style={gray!25},
  table/col sep=comma]
\addplot[omThm,very thick] table[x=inc,y=cpr]{figdata/rates-credit-risk/12-mortgage-modelling/scurve.csv};
\end{axis}
\end{tikzpicture}

\omcaption{The refinancing S-curve of the chapter's \omterm{def:rc:mortgage-modelling:model}{prepayment model}, before \omterm{def:rc:mortgage-modelling:scurve}{burnout}:
annual refinancing speed against the incentive. Turnover adds a few per cent at every
incentive. Illustrative parameters; the chapter's
tutorial.}\label{fig:rc:mortgage-modelling:scurve}
\end{omfigure}

\section{Rate paths and path dependence}

\omterm{def:rc:mortgage-modelling:scurve}{Burnout} makes each month's prepayment depend on the path so far, and the balance
left to prepay depends on all earlier prepayments: a pool cannot be valued on a tree
of rates, only by simulation of whole paths.

\begin{method}[Path-wise valuation of a pool]\label{met:rc:mortgage-modelling:paths}
Simulate monthly short-rate paths with chapter~7's \omterm{def:rc:short-rate-models:hw}{Hull--White model} (antithetic
pairs). On each path, each month: compute the ten-year par rate from the model's bond
prices, the mortgage rate, the incentive and the CPR; run the pool's amortisation
(level payment, scheduled principal, prepayment of the rest); record interest and
principal to investors. Discount each path's cash flows at its short rate plus a
spread, and average.
\end{method}

\omcode{code/firm/mbsoas/firm_mbsoas.py}{74}{94}{The pool's cash flows along every path at once.}\label{lst:rc:mortgage-modelling:flows}

\section{The option-adjusted spread in practice}

The option-adjusted spread of One Quant Book 2, chapter 21, is the constant spread
over the model's short rate that makes the average discounted cash flow equal to
the market price.

\begin{definition}[Option cost]\label{def:rc:mortgage-modelling:optioncost}
The \emph{option cost}\index{option cost} of a mortgage security is its
zero-volatility spread (the spread that reprices it on the single path of forward
rates) minus its option-adjusted spread: the value, in spread, of the borrowers'
prepayment options that volatile rates make valuable.
\end{definition}

\begin{example}[A premium pool at par]\label{ex:rc:mortgage-modelling:oas}
A pool with a WAC of 6.0\%, a coupon of 5.5\%, 348 months left, on chapter~1's SOFR curve
(ten-year par rate 3.80\% in the model's annual convention, so a mortgage rate of 5.55\% and an incentive of 0.45\%), with
Hull--White $\kappa = 3\%$, $\sigma = 90$ basis points and 4\,000 paths, priced at 100: its
option-adjusted spread is 146 basis points, its zero-volatility spread 181, and its
\omterm{def:rc:mortgage-modelling:optioncost}{option cost} 34. Its effective duration is 4.99 years, its effective convexity $-123$, its
weighted-average life 6.4 years. Without \omterm{def:rc:mortgage-modelling:scurve}{burnout} the same pool at the same spread would
be worth 99.20: faster prepayment of a premium pool returns principal at par sooner.
\end{example}

\begin{definition}[Interest-only and principal-only strips]\label{def:rc:mortgage-modelling:iopo}
An \emph{interest-only strip}\index{interest-only strip} (IO) receives a pool's interest
payments and a \emph{principal-only strip}\index{principal-only strip} (PO) its principal
payments, scheduled and prepaid; together they are the pool.
\end{definition}

\begin{example}[Splitting the pool]\label{ex:rc:mortgage-modelling:iopo}
At its option-adjusted spread the pool splits into an IO worth 25.26 and a PO worth
74.74. When rates rise 25 basis points the IO gains (26.61) and the PO loses (72.10):
slower prepayment extends the interest stream and delays the principal. The IO is a
rare asset with negative duration (\cref{fig:rc:mortgage-modelling:iopo}).
\end{example}

\begin{omfigure}
\begin{tikzpicture}
\begin{axis}[width=11.5cm,height=5.2cm,
  xlabel={parallel shift of the curve (bp)},ylabel={value (per 100 of pool)},
  tick label style={font=\small},label style={font=\small},
  xmin=-210,xmax=210,ymin=15,ymax=90,grid=major,grid style={gray!25},
  legend columns=2,legend style={font=\footnotesize,at={(0.5,-0.25)},anchor=north,draw=none,fill=none,/tikz/every even column/.append style={column sep=8pt}},
  table/col sep=comma]
\addplot[omThm,very thick,mark=*] table[x=shift,y=po]{figdata/rates-credit-risk/12-mortgage-modelling/iopo.csv};
\addplot[omDef,very thick,mark=square*] table[x=shift,y=io]{figdata/rates-credit-risk/12-mortgage-modelling/iopo.csv};
\legend{principal-only strip,interest-only strip}
\end{axis}
\end{tikzpicture}

\omcaption{Values of the pool's principal-only and \omterm{def:rc:mortgage-modelling:iopo}{interest-only strips} when the whole
curve moves, at the pool's option-adjusted spread. The PO rises steeply as rates fall
and prepayments accelerate; the IO rises with rates. Data: the chapter's
tutorial.}\label{fig:rc:mortgage-modelling:iopo}
\end{omfigure}

\section{Hedging negative convexity}

\begin{omfigure}
\begin{tikzpicture}
\begin{axis}[width=11.5cm,height=5.2cm,
  xlabel={parallel shift of the curve (bp)},ylabel={pool value},
  tick label style={font=\small},label style={font=\small},
  xmin=-210,xmax=210,ymin=87,ymax=112,grid=major,grid style={gray!25},
  legend columns=2,legend style={font=\footnotesize,at={(0.5,-0.25)},anchor=north,draw=none,fill=none,/tikz/every even column/.append style={column sep=8pt}},
  table/col sep=comma]
\addplot[omThm,very thick,mark=*,mark size=1.5pt] table[x=shift,y=value]{figdata/rates-credit-risk/12-mortgage-modelling/price.csv};
\addplot[omDef,thick,dashed] table[x=shift,y=linear]{figdata/rates-credit-risk/12-mortgage-modelling/price.csv};
\legend{pool value (OAS held),tangent at today's curve}
\end{axis}
\end{tikzpicture}

\omcaption{The pool's value against a parallel shift of the curve: below its tangent on
both sides, the signature of negative convexity. Its gains from falling rates are
capped by refinancing; its losses from rising rates grow as it lengthens. Data: the
chapter's tutorial.}\label{fig:rc:mortgage-modelling:price}
\end{omfigure}

A holder hedged to zero DV01 is short gamma: after a rise it is long duration again and
must pay fixed (or sell Treasuries) to re-hedge; after a fall it must receive. Its
hedging trades are in the direction of the move, and when many holders run the same
hedge, the trades move the market.

\begin{example}[One hundred basis points]\label{ex:rc:mortgage-modelling:event}
After a 100-basis-point rise, at the same option-adjusted spread, the pool is worth 94.50
and its duration has lengthened from 4.99 to 6.17 years. On USD~10 billion of face the
DV01 rises from USD~4.99 million to 5.83 million per basis point. A ten-year par swap on
the shifted curve has a DV01 of USD~796\,282 per billion, so the holder must pay fixed on
USD~1.06 billion of ten-year swaps to be hedged again.
\end{example}

Who runs these hedges matters as much as how large they are. A holder that does not
hedge duration, such as a central bank buying for policy reasons, absorbs negative
convexity without trading it back into the market; when that holder stops buying and
lets its holdings run off, the convexity returns to private hands, who hedge it.

\begin{omfigure}
\begin{tikzpicture}
\begin{axis}[width=11.5cm,height=5cm,
  xlabel={year},ylabel={USD trillion},
  tick label style={font=\small},label style={font=\small},
  xmin=2003,xmax=2027,ymin=0,ymax=3,grid=major,grid style={gray!25},
  xtick={2005,2010,2015,2020,2025},xticklabels={2005,2010,2015,2020,2025},
  table/col sep=comma]
\addplot[omThm,very thick] table[x=year,y=tn]{figdata/rates-credit-risk/12-mortgage-modelling/fedmbs.csv};
\end{axis}
\end{tikzpicture}

\omcaption{Mortgage-backed securities held outright by the Federal Reserve, month-end
Wednesday levels, December 2002 to August 2026: zero before 2009, a peak of USD~2.72
trillion in February 2022, USD~1.91 trillion in August 2026. Source: Federal Reserve H.4.1
via FRED (series WSHOMCB, One Quant Book 2's data file).}\label{fig:rc:mortgage-modelling:fed}
\end{omfigure}

\begin{dated}{2026-09}{Who holds the convexity}\label{dat:rc:mortgage-modelling:holders}
The Federal Reserve's holdings of agency mortgage-backed securities peaked at USD~2.72
trillion in February 2022 and stood at USD~1.91 trillion in August 2026, running off as
pools prepay. Agency mortgage securities traded about USD~367 billion a day in 2026 to
August, mostly in the to-be-announced market (SIFMA).
\end{dated}

\section{Model risk in prepayment}

Every number above rests on a behavioural model: the S-curve's centre and height, the
\omterm{def:rc:mortgage-modelling:scurve}{burnout} rate, the mortgage-rate spread. They are estimated on past refinancing waves,
and each new wave differs (technology that makes refinancing faster, tighter
underwriting that makes it slower, a \omterm{def:rc:mortgage-modelling:cc}{primary--secondary spread} that widens when lenders
are overloaded). Desks therefore report the option-adjusted spread under several
models, compute prepayment ``durations'' (the sensitivity to scaling the model's speeds),
and hold reserves for model uncertainty (chapters~26 and~27).

\section{Tutorial: a pool on many paths}\label{tut:rc:mortgage-modelling:tutorial}

\begin{tutorial}
\textbf{Goal.} Price a pass-through by option-adjusted spread, split it into strips,
and measure its negative convexity. \textbf{End state:}
\cref{fig:rc:mortgage-modelling:iopo,fig:rc:mortgage-modelling:price} and the numbers of
\cref{ex:rc:mortgage-modelling:oas,ex:rc:mortgage-modelling:event}.
\begin{tutsteps}
\item \textbf{Paths}: \texttt{simulate\_paths(curve, kappa, sigma, months, paths)} returns
monthly short rates and ten-year par rates.
\item \textbf{Flows}: \texttt{pool\_flows(pool, model, par10)}.
\item \textbf{Spread and risk}: \texttt{base()}, \texttt{duration\_at(0.01, oas)}.
\item \textbf{Hedge}: \texttt{convexity\_event()}; \texttt{fig\_rc\_mbs.py} writes the
charts.
\end{tutsteps}
\textbf{What to change next.} Double the S-curve's slope and compare the convexity;
raise the mortgage spread by 50 basis points (a widened \omterm{def:rc:mortgage-modelling:cc}{primary--secondary spread}) and
measure the change in value.
\end{tutorial}

\section{Build: the OAS engine}\label{bld:rc:mortgage-modelling:mbsoas}

\begin{build}
\bfield{Purpose} The mortgage book of the miniature firm: path-wise prepayment,
option-adjusted spreads, strips and effective risk.
\bfield{Interface} \texttt{Pool}, \texttt{PrepayModel} (with \texttt{refi},
\texttt{cpr}); \texttt{simulate\_paths}; \texttt{pool\_flows}; \texttt{pv\_paths};
\texttt{solve\_oas}; \texttt{analyse(curve, pool, model, price, kappa, sigma, paths)}.
\bfield{Rules} Monthly steps; Book~2's \texttt{firm\_prepay} for the amortisation and the
PSA ramp; common random numbers for every sensitivity; numpy only.
\bfield{Acceptance tests} \texttt{code/firm/mbsoas/tests/}: with no prepayment the pool is
an amortising loan on every path; the option-adjusted spread reprices and the strips add
up; a pool near the money has negative convexity and a positive \omterm{def:rc:mortgage-modelling:optioncost}{option cost}, its IO
gains and its PO loses when rates rise; \omterm{def:rc:mortgage-modelling:scurve}{burnout} slows prepayment.
\bfield{Stretch} A loan-level model with cohort data; a two-factor rate model; the
dollar roll's implied financing on model cash flows.
\end{build}

\begin{omsources}
\item R.\ Perli and B.\ Sack, ``Does mortgage hedging amplify movements in long-term
interest rates?'', Finance and Economics Discussion Series 2003-49, Federal Reserve Board.
\item US Department of the Treasury, \emph{Daily Treasury Par Yield Curve Rates}, 2003.
\item L.\ Hayre (ed.), \emph{Salomon Smith Barney Guide to Mortgage-Backed and
Asset-Backed Securities}, Wiley, 2001.
\item A.\ Fuster, A.\ Hizmo, L.\ Lambie-Hanson, J.\ Vickery and P.\ S.\ Willen, ``How resilient is
mortgage credit supply? Evidence from the COVID-19 pandemic'', NBER Working Paper 28843.
\end{omsources}

\section{Exercises}

\begin{exercise}[$\star$]\label{exo:rc:mortgage-modelling:1}
A pool's WAC is 6.0\% and the ten-year par rate is 3.80\%. With a mortgage spread of 1.75\%,
what is the \omterm{def:rc:mortgage-modelling:turnover}{refinancing incentive}?
\end{exercise}

\begin{exercise}[$\star$]\label{exo:rc:mortgage-modelling:2}
Using \cref{ex:rc:mortgage-modelling:oas}, what is the \omterm{def:rc:mortgage-modelling:optioncost}{option cost}, and what does it
measure?
\end{exercise}

\begin{exercise}[$\star$]\label{exo:rc:mortgage-modelling:3}
Why does the IO gain when rates rise?
\end{exercise}

\begin{exercise}[$\star\star$]\label{exo:rc:mortgage-modelling:4}
Why can a pool not be valued on a recombining tree, as a Bermudan can?
\end{exercise}

\begin{exercise}[$\star\star$]\label{exo:rc:mortgage-modelling:5}
Read \cref{fig:rc:mortgage-modelling:price}: estimate the gain from a 100-basis-point
fall and the loss from a 100-basis-point rise, and compare them.
\end{exercise}

\begin{exercise}[$\star\star$]\label{exo:rc:mortgage-modelling:6}
Without \omterm{def:rc:mortgage-modelling:scurve}{burnout} the pool is worth 99.20 at the same spread. Explain the sign.
\end{exercise}

\begin{exercise}[$\star\star\star$]\label{exo:rc:mortgage-modelling:7}
\emph{Coding.} Reproduce the re-hedge of \cref{ex:rc:mortgage-modelling:event}: the DV01
before and after the rise on USD~10 billion, and the ten-year swap notional.
\end{exercise}

\begin{exercise}[$\star\star\star$]\label{exo:rc:mortgage-modelling:8}
\emph{Find the flaw.} ``Our MBS book is duration-hedged daily, so it has no rate risk.''
\end{exercise}

\section{Problem: The Convexity Event}

\begin{problem}[{Weekend problem --- a portfolio re-hedges after a sell-off}]\label{pb:rc:mortgage-modelling:1}
A mortgage investor holds USD~10 billion of the chapter's pool, hedged to zero DV01 with
ten-year payer swaps. Rates rise by 100 basis points in parallel.

\medskip\noindent\textbf{Part I --- Before.}
\begin{enumerate}
\item Give the pool's option-adjusted spread, zero-volatility spread and \omterm{def:rc:mortgage-modelling:optioncost}{option cost}.
\item Give its duration and convexity.
\item Give the DV01 of the holding, per basis point.
\item How much of the value is interest, how much principal?
\item Which of the investor's positions is short options, and to whom?
\end{enumerate}

\medskip\noindent\textbf{Part II --- The rise.}
\begin{enumerate}[resume]
\item Give the pool's value after the rise at the same spread.
\item Give its new duration and the holding's new DV01.
\item Why did the duration lengthen?
\item Give the notional of ten-year payer swaps needed to re-hedge.
\item What happens if every mortgage investor does the same trade on the same day?
\end{enumerate}

\medskip\noindent\textbf{Part III --- Strips.}
\begin{enumerate}[resume]
\item Give the IO's and PO's values at $\pm25$ basis points.
\item Which strip hedges the other's convexity?
\item Why do IOs appeal to banks with deposits whose value rises with rates?
\item What does a PO holder lose if refinancing slows for reasons unrelated to rates?
\item Which strip is more exposed to prepayment-model risk?
\end{enumerate}

\medskip\noindent\textbf{Part IV --- Judgement.}
\begin{enumerate}[resume]
\item Should the investor hedge with swaptions instead of swaps? What would it cost?
\item Why did the 2003 episode last months rather than days?
\item How would a central bank holding mortgage securities change the dynamics?
\item State the \emph{named result}: the DV01 before and after and the swap notional of
the re-hedge.
\item In one sentence: why does a mortgage lengthen when you least want it to?
\end{enumerate}
\end{problem}

\section{Interview questions}

\begin{interviewq}[$\star$ \iqroles{trader, researcher}]\label{iq:rc:mortgage-modelling:1}
Why do mortgage-backed securities have negative convexity?
\end{interviewq}

\begin{interviewq}[$\star\star$ \iqroles{researcher}]\label{iq:rc:mortgage-modelling:2}
What goes into a \omterm{def:rc:mortgage-modelling:model}{prepayment model}, and which parts depend on rates?
\end{interviewq}

\begin{interviewq}[$\star\star$ \iqroles{researcher, developer}]\label{iq:rc:mortgage-modelling:3}
How do you compute an option-adjusted spread, and why Monte Carlo?
\end{interviewq}

\begin{interviewq}[$\star\star$ \iqroles{trader}]\label{iq:rc:mortgage-modelling:4}
Rates sell off 50 basis points. What do mortgage hedgers do, and what does it do to the
market?
\end{interviewq}

\begin{interviewq}[$\star\star\star$ \iqroles{risk, researcher}]\label{iq:rc:mortgage-modelling:5}
How would you measure and reserve for prepayment-model risk?
\end{interviewq}

\begin{interviewq}[$\star\star\star$ \iqroles{developer}]\label{iq:rc:mortgage-modelling:6}
Your effective duration jumps from day to day by half a year with no market move. What
do you check?
\end{interviewq}
